\section{Marco referencial}

\subsection{Marco histórico}

La preocupación por adaptar el trabajo y los entornos de estudio a las capacidades del cuerpo humano no es nueva. El término \emph{ergonomía} fue acuñado en 1857 por el polaco Wojciech Jastrzębowski, quien propuso estudiar el trabajo como una ciencia, y solo fue retomado en 1949, cuando Murrell lo reintrodujo y se fundó en el Reino Unido la primera sociedad dedicada a la disciplina, la Ergonomics Research Society [6]. Desde entonces, la ergonomía se consolidó como un campo científico que estudia cómo las condiciones físicas de una actividad influyen en la salud de quien la realiza.

Con la llegada de la era digital, el computador pasó a ser una herramienta central de trabajo y de estudio, y con ello aparecieron problemas de salud asociados a su uso prolongado. La Asociación Americana de Optometristas definió el síndrome visual informático como el conjunto de problemas visuales y oculares relacionados con el uso prolongado de pantallas, cuyos signos más frecuentes son fatiga visual, cefalea, visión borrosa, ojos secos y dolor de cuello y hombros [7]. Así, a la preocupación clásica por la postura se sumó la exposición a pantallas como factor de riesgo.

Un cambio reciente en la escala de esta exposición fue la educación virtual durante la pandemia de la COVID-19, que llevó a los estudiantes universitarios a pasar la mayor parte de su jornada académica frente a dispositivos digitales. Estudios de ese periodo documentaron trastornos musculoesqueléticos y malas prácticas ergonómicas en universitarios [1], así como una alta presencia del síndrome visual informático en estudiantes de posgrado y de medicina [2], [3].

Este recorrido ubica el punto de partida del estado del arte: la relación entre exposición a pantallas, postura y síntomas físicos en estudiantes universitarios, tal como se estudia desde la virtualización de la educación hasta hoy.

\subsection{Estado del arte}
\estado{En construcción}
\notafase{Tema 4.2}

\subsection{Marco teórico}
\estado{En construcción}
\notafase{Tema 4.3}

\subsection{Marco conceptual}
\estado{En construcción}
\notafase{Tema 4.4}

\subsection{Marcos complementarios}
\estado{En construcción}
\notafase{Tema 4.5}
